\section{Conclusiones y Trabajos Futuros}

\subsection{Resumen de Contribuciones}

Resulta revelador observar cómo una integración pensada con cuidado entre redes terrestres y no terrestres puede transformar las capacidades de monitoreo en Smart Cities. Este trabajo ha propuesto y validado una arquitectura híbrida TN-NTN que logra cobertura efectiva superior al 97\%, reduce latencia promedio de forma significativa y mejora la eficiencia energética global del sistema IoT en un 23\% respecto a soluciones puramente terrestres bajo condiciones realistas de carga.

Las contribuciones centrales incluyen un modelo de sistema multi-capa con orquestación contextual inteligente, mecanismos de handover make-before-break asistidos por predicción, y un framework de evaluación que combina simulaciones de alta fidelidad 3D con métricas orientadas a aplicaciones urbanas. Más allá de los números, la propuesta demuestra que es posible conciliar cobertura ubicua, resiliencia y eficiencia operativa sin caer en soluciones excesivamente complejas o costosas. Los resultados no solo confirman las ventajas teóricas de NTN, sino que las aterrizan en escenarios de monitoreo práctico con miles de dispositivos.

\subsection{Recomendaciones}

Uno de los desafíos persistentes al pasar de la simulación al despliegue radica en la brecha entre modelos ideales y condiciones operativas reales. Para operadores y autoridades urbanas interesadas en adoptar enfoques similares, recomendamos comenzar con pilotos a escala media en distritos representativos, priorizando aplicaciones críticas de resiliencia y monitoreo ambiental.

Carbonara et al. \cite{Carbonara2025} enfatizan la importancia de plataformas de testing integradas TN-NTN; coincidimos plenamente y sugerimos emplear entornos de prueba híbridos que repliquen tanto movilidad vehicular como patrones de tráfico IoT bursty. Desde el punto de vista de implementación, resulta aconsejable adoptar un enfoque incremental: primero habilitar multi-conectividad básica, luego incorporar orquestación basada en reglas y, finalmente, evolucionar hacia agentes de decisión más avanzados. 

La selección cuidadosa de constelaciones (combinando LEO para rendimiento y GEO para respaldo) junto con políticas energéticas estrictas en los dispositivos se revela clave para mantener viabilidad económica y operativa a largo plazo. No subestimen tampoco los aspectos regulatorios y de espectro, que en muchos casos estudiados condicionan más el éxito que los propios desafíos técnicos.

\subsection{Direcciones Futuras}

Lejos de ser un asunto resuelto, la convergencia entre NTN e IoT abre un amplio horizonte de investigación. Aunque nuestra arquitectura marca un avance concreto, persisten oportunidades significativas para refinamiento y extensión.

Resulta particularmente prometedor explorar la integración nativa con tecnologías 6G emergentes como Reconfigurable Intelligent Surfaces (RIS) para mejorar enlaces satelitales en entornos densamente obstruidos, o el uso de aprendizaje federado para optimizar decisiones de handover de forma distribuida preservando privacidad. La escalabilidad a decenas de miles de dispositivos por km² bajo patrones de tráfico reales también merece atención prioritaria.

Ogbodo et al. \cite{Ogbodo2022} ofrecen una visión amplia sobre LPWAN e IoT en Smart Cities que sigue siendo relevante; alineamos nuestros trabajos futuros con esa perspectiva al proponer estudios que incorporen mayor heterogeneidad de dispositivos, integración con edge computing para procesamiento local de datos y evaluaciones económicas completas de TCO (Total Cost of Ownership) en diferentes contextos urbanos.

Otras líneas interesantes incluyen la robustez ante amenazas de seguridad específicas de NTN, la optimización conjunta de energía harvesting y protocolos de transmisión, y la validación en campo a gran escala. Creemos firmemente que la verdadera inteligencia de las ciudades del futuro no residirá solo en la densidad de sensores, sino en su capacidad de permanecer conectados y útiles bajo cualquier circunstancia. Esta investigación representa un paso en esa dirección, pero el camino aún ofrece numerosos desafíos estimulantes y oportunidades de impacto real.